\problem{Стабильность звезды}{10.0}

\part{Параллакс}

\begin{answerbox}{2.5cm}
Формула: \(d=\)

\bigskip

Численное значение: \(d=\)
\end{answerbox}

\part{Формирование протозвезды}

\begin{answerbox}{1.5cm}
\begin{multicols}{2}
\(N=\)

\columnbreak

\(U=\)
\end{multicols}
\end{answerbox}

\begin{answerbox}{1.5cm}
\(g(r)=\)
\end{answerbox}

\begin{answerbox}{1.5cm}
\(f=\)
\end{answerbox}

\begin{answerbox}{1.5cm}
\(M_J=\)
\end{answerbox}

\part{Параметры Солнца}

\begin{answerbox}{1.7cm}
\(\dfrac{dP}{dr}=\)
\end{answerbox}

\begin{answerbox}{1.5cm}
\begin{multicols}{2}
\(E_G=\)

\columnbreak

\(\beta=\)
\end{multicols}
\end{answerbox}

\begin{answerbox}{2.5cm}
Формула: \(\langle P\rangle=\)

\bigskip

Численное значение: \(\langle P\rangle=\)
\end{answerbox}

\begin{answerbox}{1.5cm}
\(T_0=\)
\end{answerbox}

\part{Излучение Солнца}

\begin{answerbox}{1.5cm}
\(T_s=\)
\end{answerbox}

\begin{answerbox}{1.5cm}
\begin{multicols}{2}
\(\vec D=\)

\columnbreak

\(\langle\vec D\rangle=\)
\end{multicols}
\end{answerbox}

\begin{answerbox}{1.5cm}
\begin{multicols}{2}
\(\langle D^2\rangle=\)

\columnbreak

\(t=\)
\end{multicols}
\end{answerbox}

\begin{answerbox}{1.5cm}
\begin{multicols}{2}
\(l=\)

\columnbreak

\(t=\)
\end{multicols}
\end{answerbox}

\begin{answerbox}{1.5cm}
\begin{multicols}{2}
\(C=\)

\columnbreak

\(n=\)
\end{multicols}
\end{answerbox}

\begin{answerbox}{1.5cm}
На основании графика: \(n=\)
\end{answerbox}
